\section{Local work under a neighborhood-type promise}
\label{sec:op3-neighborhood-types}

\paragraph{Source definition and transfer scope.}
For an uncolored graph, two vertices are of the same neighborhood type
when $N(u)\setminus\{v\}=N(v)\setminus\{u\}$.
Definitions 1--2 of \citet{lampis2009types}, PDF p.\ 4, define neighborhood
diversity by a partition into such types. The source's Theorem 5,
pp.\ 12--13, describes clique or independent types and complete or empty
connections between them. Its global partition-computation and logical
model-checking algorithms are not local PageRank algorithms and are not
used in the following work bound. Here the type partition is a graph
promise, not input or free preprocessing.

\begin{theorem}[Unsupplied neighborhood types; \statusproved]
\label{thm:op3-local-neighborhood-types}
Suppose the original connected simple unweighted graph has a partition
into $k$ neighborhood types. The producer of
\cref{thm:op3-frontier-response-groups}, given only its original local
graph interface, seed, $0<\alpha\leq1$ and $0<e<1$, returns original
ACL output with $V<2/e$ and work and total allocated words
\[
 O\bigl((1+V)(1+k)^2\log(2+V)\bigr)
 \subseteq O\bigl((1+e^{-1})(1+k)^2\log(2+e^{-1})\bigr).
\]
The producer need not know $k$ or any type membership. All graph queries,
group refinement, transformed-load reporting, history, reconstruction and
output are charged. The bound is deterministic and independent of target
teleportation and ambient size in the exact-real word model.
\end{theorem}
\begin{proof}
The initial seed group is the only known group. If it is ineligible, the
producer stops after a degree query. Otherwise the seed is admitted first.
Thereafter, for each promised global type, its unadmitted members are either
all undiscovered or all known in one response group. Indeed, the first
admitted vertex adjacent to any member of a still-undiscovered type is
adjacent to every other member: the common-type equality applies outside
the pair. Its original row therefore discovers them together. If the
admitted vertex itself belongs to that type, the same statement concerns
the remaining members; a clique type exposes all of them and an independent
type exposes none. The seed has already left the frontier, so its distinct
source load causes no exception among these remaining members.

Inductively, a later admitted row treats every still-unadmitted member of
a global type identically. Hence the row-driven split either moves all of
that type's remaining members to the same child or moves none. A type can
never be divided between two completed-state response groups. Several
global types may share one group, but each nonempty group contains a type
that occurs in no other group. There are therefore at most $k$ nonempty
groups after each completed pivot. A pivot starts with at most $k$ groups,
creates at most one child per old group and at most one fresh discovery
group. Before retiring empty parents its live list has size at most
$2k+1$. Thus every $b_i$ in
\cref{thm:op3-frontier-response-groups} is at most $2k+1$.
Its proved volume, work and allocation bounds give the result. This
argument uses the promised partition only as a comparison object; the
producer's code never reads it.
\end{proof}

\begin{proposition}[Four types retain the output lower bound; \statusproved]
\label{prop:op3-four-type-output-lower-bound}
For every $0<e\leq1/16$ and $0<\alpha\leq1/3$, there are original
graphs with at most four neighborhood types and arbitrarily large ambient
volume on which every nonnegative original ACL output must contain at least
$1/(16e)$ positive coordinates. Consequently the inverse-accuracy power in
\cref{thm:op3-local-neighborhood-types} is optimal, up to logarithmic
factors, already for fixed $k=4$ in this parameter range.
\end{proposition}
\begin{proof}
Let $m=\lfloor1/(8e)\rfloor\geq1/(16e)$. Join a seed center $v$ to
$m$ private leaves and to another center $h$, and join $h$ to an arbitrary
number $D\geq1$ of other private leaves. The two centers and their two
leaf banks form four neighborhood types. Increasing $D$ increases ambient
volume without changing $d_v=m+1$.

For any $x\geq0$ with $0\leq e_v-(D_G-\gamma A)x\leq e d$, the seed
inequality gives $x_v\geq1/d_v-e$. Here $D_G$ denotes the original degree
matrix, distinguished from the remote leaf count. Since
\[
 e d_v\leq\tfrac18+e\leq\tfrac3{16}
 <\tfrac13\leq\frac\gamma{1+\gamma},
\]
an omitted near leaf would have residual
$\gamma x_v\geq\gamma(1/d_v-e)>e$, a contradiction. All $m$ near
leaves are mandatory. This extends the same original-seed inequality used
in \cref{prop:op3-large-ambient-leaf-lower-bound}, while keeping a constant
number of global types. It is an output-writing lower bound, not a claim
that every instance requires that much work.
\end{proof}

\paragraph{Implemented audit; Measured.}
\path{frontier_neighborhood_types.py} generates 349 graphs from all
connected quotient shapes on at most four vertices, clique/independent
type choices and recorded multiplicity patterns. It checks every physical
source at the listed parameters, totaling 20,826 original ACL outputs and
157,500 original residual rows. The partition is validator-only information.
There are 71,050 group-count checks and 194,276 checks that unadmitted types
remain unsplit and are discovered together. The run also records 12,959
states with unequal original degrees inside one response group. Twenty-one
selected cases have independent dense obstacle/Schur validation and 164
reverse coordinate/aggregate checks.

Private star and double-star oracles test huge ambient leaf banks without
reading the remote hub row. The double-star tests also verify the mandatory
near-leaf inequality and every potentially nonzero original residual;
unseen remote leaves have no active original neighbor and hence zero
residual. Parameters, configurations, counters and hashes are in
\path{FRONTIER_NEIGHBORHOOD_TYPES_AUDIT.json}. This is a structural upper
bound with a matching accuracy power. General OP3 remains \statusopen.
